\textbf{Lessons for teams building tool-enabled agents.}

\emph{L1: never execute the retriever's top hit.} Retriever-only routing, the cheapest
design and a tempting one for latency (a guard on it still needs an intent call), escalates side effects on \GRNonePoolConf\,\% of requests (verified) and selects a writing tool for
\RWRetRate\,\% of read requests. An
LLM selector over a short list lowers that risk but does not remove it (Table~\ref{tab:main}),
and the escalations that remain come from both stages: a selector that picks the stronger sibling
with the right tool in view, and a retriever that leaves the right tool out while the selector
rarely declines (Table~\ref{tab:split}).

\emph{L2: verify side-effect metrics before trusting them, with more than one rater.} Scored
automatically, escalation appeared to grow with catalog size, and two of every three flags were
false. Severity labels inherit the ecosystem's noisy annotations, so an evaluation of tool search
should report human-verified escalations, or at least the precision of its severity labels. Humans
disagree too: our two raters agreed only slightly ($\kappa=\IndKappa$) before a joint review resolved all but one
case, and one rater's blind
labels would have multiplied the sensitivity estimate several-fold. Teams should write the rubric down,
double-rate a sample and report how much the result depends on the rater.

\emph{L3: put side-effect metadata in the policy, not in the prompt.} The servers' own
annotations and our inferred tags moved accuracy by about a point and did not measurably change
escalation: the verb in the tool name already carries that information. Metadata pays off where code acts on it, in the guard.

\emph{L4: lengthen the shortlist with siblings, and test it on your model.} Sibling expansion gave
2.6--4.5 accuracy points for the three larger selectors. A plain shortlist of the same length gave
Gemma3 and Phi-4 nothing or less, and recovered most of the gain only for Qwen3.5, so ``just
retrieve more'' is not a reliable substitute. Siblings are what the selector must tell apart, which
makes them a natural way to fill a longer list, but we have not isolated their effect from that of
the tags and ordering the sibling variant also adds.
Neither change moves escalation measurably.

\emph{L5: choose the guard by the confirmations you can afford.} Holding every increase in side
effect removes almost all escalation but asks for confirmation on 15--20\,\% of calls, depending on the selector; guards
that only police reads or deletions ask less often but miss the most common real case, a read
request answered by a create. Most needless holds trace to fuzzy operation boundaries, not to the
guard itself. Table~\ref{tab:checklist} condenses the lessons into a checklist.

\begin{table}[t]\centering
\caption{Checklist for teams building tool-selecting agents, from our results.}\label{tab:checklist}
\scriptsize\setlength{\tabcolsep}{3pt}
\begin{tabular}{p{0.36\columnwidth}p{0.58\columnwidth}}\toprule
situation & what our evidence supports\\\midrule
tools chosen by embedding similarity & do not execute the top hit unchecked; add an LLM selector and a guard (L1)\\
tools carry read-only or destructive hints & enforce them in code, never trust them blindly; many are missing or wrong (L3)\\
similar operations on one resource & compare sibling-aware shortlists with a plain shortlist of the same length on your model (L4)\\
a stronger side effect is possible & choose a guard policy and measure how many correct calls it holds (L5)\\
evaluating tool selection & report read-to-write violations, human-verified, with how much they depend on the rater (L2)\\
\bottomrule\end{tabular}\end{table}

\textbf{Where this applies.} The consequences are largest for cloud and infrastructure consoles,
source hosting and CI automation, issue trackers and messaging. The catalog contains such tools,
although authenticated servers of these kinds are under-represented (Section~\ref{sec:threats}). An operator who asks an agent to inspect a deployment or list open
incidents should not have the request routed to the tool that restarts or closes them; in our benchmark a
retriever-only router selects such a tool for \RWRetRate\,\% of read requests (verified). These are
selections on a targeted benchmark; we executed no tool and observed no production incident.

\textbf{Implications for tool and protocol designers.} Three ecosystem habits make side effects
hard to police. (i)~Annotations are optional and often absent, including on tools that
delete. (ii)~\texttt{destructiveHint} conflates ``overwrites an existing value'' with ``deletes
or cannot be undone'', so a guard keyed on it holds ordinary edits. (iii)~\emph{Mode-switch} tools
(a cancel tool with a preview flag, an \texttt{action} parameter selecting list, set or delete)
give one name to operations of different severity, so neither retrieval nor a guard can tell a
read from a delete before the arguments are filled. Splitting such tools and adding an explicit
operation class to tool metadata would let hosts enforce what the hints only suggest.
